\chapter{Problem Formulation}
\label{ch:problem_formulation}

This chapter establishes the biological foundations and control engineering setting of the thesis.
It introduces neural activity and its representation across scales, explains the recruitment of brain regions during focal seizures and describes the available measurement and stimulation interfaces.
The final section places these elements in the simulation and feedback control setting used in this work.

\section{Neural Activity and Modeling Across Scales}
\label{sec:neural_activity_background}

Neural activity comprises the electrical signaling of individual neurons and the collective responses of their interacting populations.
A neuron maintains an electrical potential difference across its cell membrane.
Changes in ionic currents across that membrane alter the membrane potential and can generate an action potential, a brief electrical pulse that travels along the axon~\cite{Hodgkin1952}.
At a chemical synapse, activity in the sending neuron changes the input received by another neuron.
Excitatory inputs tend to promote firing, whereas inhibitory inputs tend to restrain it~\cite{Wilson1972}.
At this cellular scale, conductance-based descriptions such as the Hodgkin-Huxley model describe how ionic currents generate these signals~\cite{Hodgkin1952}.

Neurons participate in circuits with recurrent excitation and inhibition.
Their collective activity can exhibit oscillations even though individual neurons need not fire on every oscillation cycle~\cite{Buzsaki2012}.
Neural mass models describe these collective dynamics through population quantities such as mean firing rate and mean postsynaptic potential~\cite{SanzLeon2013}.
A postsynaptic potential is the membrane-potential response to synaptic input.
Population descriptions retain the interactions between groups of neurons while leaving individual spike timings unresolved.
Wilson-Cowan models describe interacting excitatory and inhibitory population activity, while Jansen-Rit models include synaptic filtering and population membrane potentials~\cite{Wilson1972,Jansen1995}.

In cortical neural mass models, pyramidal cells form an important excitatory population that interacts with local interneurons~\cite{Jansen1995}.
A cortical macro-column is represented as a local circuit of such populations extending through the layers of the cortex, the folded outer layer of the cerebrum.
Using a neural mass to describe an entire brain region is a further coarse-graining: its variables represent regional activity aggregated over many local circuits~\cite{SanzLeon2013,SanzLeon2015}.
Figure~\ref{fig:neural_scales} illustrates this relationship between regional network nodes, interacting populations and their constituent neurons.

\begin{figure}
    \centering
    \tikzsetnextfilename{problem_overview}
\begin{tikzpicture}[
    thesis figure,
    zoom/.style={draw=accentcolor!65!black, dashed, line width=0.7pt},
    region/.style={circle, draw=black!65, fill=black!25,
        minimum size=4mm, inner sep=0pt},
    pics/neuron/.style={code={
        \draw[accentcolor!85!black, line width=0.6pt]
            (0,0) circle (0.065)
            (0.065,0) -- (0.42,0) -- (0.52,0.10)
            (0.42,0) -- (0.52,-0.10)
            (-0.05,0.04) -- (-0.18,0.16) -- (-0.28,0.15)
            (-0.18,0.16) -- (-0.17,0.25)
            (-0.06,-0.02) -- (-0.20,-0.10) -- (-0.29,-0.06)
            (-0.20,-0.10) -- (-0.22,-0.20);
    }}
]
    \node at (1.65,1.75) {Whole-brain network};
    \node at (6.05,1.75) {Regional populations};
    \node at (10.85,1.75) {Constituent neurons};

    \foreach \side in {-1,1} {
        \begin{scope}[shift={(1.65,0)}, xscale=\side]
            \path[draw=black!40, fill=black!3]
                (0.05,1.25)
                .. controls (0.45,1.57) and (1.10,1.20) .. (1.30,0.65)
                .. controls (1.65,0.05) and (1.20,-1.20) .. (0.50,-1.33)
                .. controls (0.13,-1.40) and (0.03,-0.95) .. (0.05,-0.65)
                -- cycle;
        \end{scope}
    }
    \foreach \name/\x/\y in {a/0.70/0.65,b/1.32/0.90,c/2.30/0.80,d/2.48/0.30,e/0.65/-0.25,f/1.35/-0.65,g/2.40/-0.65} {
        \coordinate (\name) at (\x,\y);
    }
    \foreach \a/\b in {a/b,a/e,a/f,b/c,b/f,c/d,c/g,d/f,d/g,e/f,f/g} {
        \draw[black!40] (\a) -- (\b);
    }
    \foreach \name in {a,b,c,e,f,g} {
        \node[region] at (\name) {};
    }
    \node[region, fill=accentcolor!75, draw=accentcolor!85!black,
        minimum size=4.5mm] at (d) {};
    \draw[zoom] (d) circle (0.27);

    \draw[black!40, fill=black!3] (6.05,0) circle (1.45);
    \draw[->, accentcolor!70!black, line width=0.8pt]
        (5.75,0.05) to[bend right=20] (6.65,-0.48);
    \draw[-{Bar}, orange!80!black, line width=0.8pt]
        (6.72,-0.35) to[bend right=25] (5.95,0.43);
    \draw[->, accentcolor!70!black, line width=0.8pt]
        (5.30,-0.65) -- (5.52,0.03);
    \foreach \dx/\dy in {-0.14/0.10,0.12/0.10,-0.14/-0.12,0.12/-0.12} {
        \fill[accentcolor!75] (5.68+\dx,0.35+\dy) circle (0.065);
        \fill[orange!80!black] (6.80+\dx,-0.55+\dy) circle (0.065);
        \fill[accentcolor!75] (5.25+\dx,-0.85+\dy) circle (0.065);
    }
    \draw[zoom] (5.68,0.35) circle (0.43);

    \draw[black!40, rounded corners=5pt, fill=black!3]
        (9.15,-1.30) rectangle (12.55,1.30);
    % Zoom connectors link the selected node and population to their expansions.
    \draw[zoom] (2.66,0.47) -- (4.82,0.90);
    \draw[zoom] (2.66,0.13) -- (4.82,-0.90);
    \draw[zoom] (6.08,0.51) -- (9.15,1.30);
    \draw[zoom] (6.08,0.19) -- (9.15,-1.30);
    \foreach \x in {9.65,10.65,11.65} {
        \foreach \y in {-0.78,0,0.78} {
            \pic at (\x,\y) {neuron};
        }
    }
    \node[thesis secondary, align=center] at (1.65,-1.95)
        {Connected regions};
    \node[thesis secondary, align=center] at (6.05,-1.95)
        {Interacting excitation\\and inhibition};
    \node[thesis secondary, align=center] at (10.85,-1.95)
        {Many cells within\\one population};
\end{tikzpicture}

    \caption{Neural activity across scales: Zoom connectors expand a selected node of a whole-brain network into interacting regional populations, then expand one population into its constituent neurons. Neural mass models summarize population activity and regional nodes aggregate the dynamics of local circuits.}
    \label{fig:neural_scales}
\end{figure}

At the whole-brain scale, regions communicate through bundles of axons, including long-range white-matter pathways~\cite{Hagmann2008}.
Whole-brain network models connect regional models through a structural connectome, a representation of anatomical connections between brain regions~\cite{SanzLeon2013}.
Constructing this connectome involves a parcellation that divides the brain into anatomical units and an estimate of the pathways connecting them.
Diffusion magnetic resonance imaging measures direction-dependent water diffusion in tissue, from which local fiber orientations can be estimated.
Tractography follows these estimated orientations to reconstruct candidate white-matter pathways.
Combining these pathways with the parcellation yields regional connection weights and estimates of pathway lengths~\cite{Hagmann2008}.

In a dynamical network, connection weights determine the strength of interactions, while pathway lengths and conduction speeds determine transmission delays~\cite{SanzLeon2015}.
Such networks have been used to investigate how regional excitability and anatomical connectivity shape seizure recruitment~\cite{Olmi2019}.

\section{Focal Seizures and Recruitment}
\label{sec:seizure_recruitment_background}

The regional interactions described above also provide pathways for pathological activity to spread.
An epileptic seizure is a temporary episode of clinical signs or symptoms arising from abnormally excessive or synchronized neuronal activity~\cite{Fisher2005}.
The term focal describes seizures that originate within networks limited to one hemisphere; the activity may subsequently spread beyond its initial territory~\cite{Fisher2017}.
The resulting manifestations depend on the neural systems involved~\cite{Fisher2017}.

For the network description adopted here, the \ac{ez} denotes the regions that initiate seizure activity.
Recruitment means that a region outside the initiating zone develops seizure-like activity through its interaction with the active network.
It is a change in that region's dynamics, rather than merely the detection of an electrical signal generated elsewhere~\cite{Olmi2019}. % what is "it"? the version before was clearer

Whether a region is recruited depends on both its local susceptibility and the input it receives from other regions.
Structural connectivity alone therefore does not determine the propagation pattern~\cite{Olmi2019}.
Population models can represent differences in susceptibility through regional parameters controlling the balance of excitation and inhibition~\cite{Wendling2002,Yu2024}.
This combination of local dynamics and network interactions makes propagation a dynamical problem.

The primary control objective is to keep recruitment outside the \ac{ez} as low as possible.
Complete prevention is the ideal outcome, while a reduction in recruitment also represents improvement relative to the uncontrolled case.
The secondary objective is minimal intervention, expressed through low stimulation energy and small overall electrode-current magnitudes.

\section{Measurement and Stimulation}
\label{sec:measurement_intervention_background}

Acting on recruitment requires a way to observe neural activity and an intervention that can influence its evolution.
\Ac{eeg} records electrical potential differences between electrodes on the scalp.
The recorded potentials reflect the superposition of contributions from neural currents after conduction through the tissues of the head~\cite{Buzsaki2012}.
Consequently, an electrode does not directly measure the activity of a single underlying region.
The relationship between regional recruitment and the available measurements is therefore indirect~\cite{Buzsaki2012}.

\Ac{tes} provides the intervention route by applying currents through electrodes on the scalp.
These currents produce electric fields within the head, with their spatial distribution depending on the electrode arrangement and the electrical properties of the tissues~\cite{Huang2019}.
The effects of stimulation on collective neural activity can be investigated by combining a head model with a model of neural populations~\cite{Merlet2013,Kunze2016}.
Admissible electrode currents must additionally satisfy the imposed stimulation limits.

\section{Control Engineering Setting}
\label{sec:control_engineering_setting}

The measurement and stimulation interfaces define the information and intervention available for feedback control.
This thesis studies that setting in simulation on a biophysical plant model comprising a network of regional neural masses, a stimulation mapping and an \ac{eeg} measurement mapping.
The plant represents the system on which stimulation acts and against which control performance is evaluated.
Its regional activity is available for assessing recruitment, while the controller receives electrical measurements.

The feedback controller uses \ac{mpc} to select electrode currents by predicting the effects of candidate stimulation sequences.
At each update, \ac{mpc} evaluates these sequences over a finite horizon according to a control objective and constraints~\cite{Rawlings2017}.
A healthy reference specifies the desired non-seizure activity against which predicted measurements are assessed.
The control objective also accounts for stimulation expenditure, while constraints restrict the admissible electrode currents.
Only the first input of the selected sequence is applied and the optimization is repeated after new measurements become available~\cite{Rawlings2017}.

To obtain a prediction model for the \ac{mpc}, simulations of the plant generate data consisting of applied currents and the resulting \ac{eeg} measurements.
System identification uses input-output data to determine a dynamical model~\cite{Ljung1999}.
The identified prediction model forecasts future measurements under candidate stimulation currents.
It is distinct from the biophysical plant; the plant's neural states and mechanistic equations are not supplied to the controller.
Identification takes place offline, before closed-loop evaluation.
During feedback control, measurement processing prepares the incoming \ac{eeg} for the prediction model.

Recruitment is assessed from regional activity in the plant, whereas control decisions rely on the available measurements and the identified prediction model.
Figure~\ref{fig:control_setting} shows the resulting feedback loop and the offline identification procedure.
The subsequent chapters specify the biophysical plant, prediction model identification and the predictive control formulation.

\begin{figure}
    \centering
    \begin{tikzpicture}[
    >=Stealth,
    thesis figure,
    block/.style={draw=black!70, fill=white, align=center,
        minimum height=0.9cm, inner sep=4pt},
    signal/.style={->, line width=0.8pt, draw=black!85},
    offline/.style={->, dashed, line width=0.8pt, draw=accentcolor!85!black},
    group/.style={draw=black!45, rounded corners=3pt, fill=black!3}
]
    \node[anchor=west, thesis heading] at (0,0.45) {Online feedback};

    \path[group, fill=accentcolor!4] (1.6,0) rectangle (6.9,-5.25);
    \node[anchor=west] at (1.8,-0.3) {\acs{mpc}};
    \node[block, minimum width=2.0cm] (objective) at (3.1,-1.15)
        {Objective};
    \node[block, minimum width=2.0cm] (constraints) at (5.4,-1.15)
        {Constraints};
    \node[block, minimum width=4.5cm] (optimize) at (4.25,-2.7)
        {Optimization};
    \node[block, minimum width=4.5cm] (prediction) at (4.25,-4.55)
        {Prediction model};

    \node[align=center, thesis secondary] at (0.75,-0.65)
        {Healthy\\reference};
    \draw[signal] (0,-1.15) -- (objective.west);
    \draw[signal] (objective.south) -- (3.1,-2.25);
    \draw[signal] (constraints.south) -- (5.4,-2.25);
    \draw[signal] (2.4,-3.15) -- (2.4,-4.10);
    \node[thesis secondary, align=center] at (3.2,-3.63)
        {Candidate\\currents};
    \draw[signal] (6.1,-4.10) -- (6.1,-3.15);
    \node[thesis secondary, align=center] at (5.1,-3.63)
        {Predicted\\measurements};

    \path[group] (8.2,-1.65) rectangle (15.7,-3.75);
    \node[anchor=west, thesis secondary] at (8.4,-1.90)
        {Biophysical plant model};
    \node[block, minimum width=2.0cm] (stim) at (9.4,-2.7)
        {Stimulation\\mapping};
    \node[block, minimum width=2.1cm] (neural) at (11.95,-2.7)
        {Regional\\neural dynamics};
    \node[block, minimum width=2.0cm] (measure) at (14.5,-2.7)
        {\acs{eeg}\\measurement};
    \draw[signal] (optimize) -- (stim);
    \node[thesis secondary] at (7.55,-2.4) {First input};
    \draw[signal] (stim) -- (neural);
    \draw[signal] (neural) -- (measure);

    \node[block, minimum width=3.6cm] (process) at (11.7,-4.55)
        {Measurement processing};
    \draw[signal] (measure.east) -- (15.9,-2.7) |- (process.east);
    \node[thesis secondary, anchor=west] at (14.6,-4.25) {\acs{eeg}};
    \draw[signal] (process.west) -- (prediction.east);

    \draw[black!25] (0,-5.85) -- (16.2,-5.85);
    \node[anchor=west, thesis heading] at (0,-6.25)
        {Offline identification};
    \node[block, minimum width=4.5cm] (identify) at (4.25,-7.15){Training};
    \node[block, minimum width=4.5cm] (data) at (9.4,-7.15)
        {Dataset generation};
    \draw[offline, preaction={draw=white, line width=2.5pt}]
        (9.4,-3.75) -- (data.north);
    \draw[offline] (data) -- (identify);
    \draw[offline] (identify.north) -- (prediction.south);
    \node[thesis secondary, text=accentcolor!85!black, anchor=west]
        at (4.4,-5.58) {Parameters};
\end{tikzpicture}

    \caption{Control engineering setting: Solid arrows show online signal flow. Within \acs{mpc}, the control objective and constraints guide optimization using the prediction model; the healthy reference enters the control objective from outside the controller. The first input of the selected sequence is applied to the biophysical plant and processed \acs{eeg} measurements return to the prediction model. Below the feedback loop, dashed arrows show offline data generation and system identification.}
    \label{fig:control_setting}
\end{figure}
